% Part: incompleteness
% Chapter: incompleteness-provability
% Section: lob-thm

\documentclass[../../../include/open-logic-section]{subfiles}

\begin{document}

\olfileid{inc}{inp}{lob}

\olsection{Teorema L\"ob}

Ukara G\"odel kanggo teori~$\Th{T}$ iku titik tetep saka $\lnot
\OProv[\Th{T}](y)$, yaiku !!a{sentence}~$!G$ kang nyukupi
\[
\Th{T} \Proves \lnot \OProv[\Th{T}](\gn{!G}) \liff !G.
\]
Yen teori iki konsisten, ukara kasebut ora !!{derivable}, amarga
yen $\Th{T} \Proves !G$, (a) miturut
sarat !!{derivability}~(1), $\Th{T} \Proves \OProv[\Th{T}](\gn{!G})$,
lan (b) saka $\Th{T} \Proves !G$ bebarengan karo $\Th{T} \Proves
\lnot \OProv[\Th{T}](\gn{!G}) \liff !G$, kita oleh $\Th{T} \Proves
\lnot \OProv[\Th{T}](\gn{!G})$; mula $\Th{T}$ mesthi ora konsisten. Cathetan koreksi sumber OLPL-800: ukara sumber nyebut G ora bisa dibuktekake tanpa syarat, nanging argumentasine mung ngasilake kontradiksi yen T konsisten. Terjemahan iki nyebut syarat konsistensi kanthi cetha.
Saiki lumrah yen kita takon kepriye status titik tetep saka
$\OProv[\Th{T}](y)$, yaiku !!a{sentence}~$!H$ kang nyukupi
\[
\Th{T} \Proves \OProv[\Th{T}](\gn{!H}) \liff !H.
\]
Yen ukara iki !!{derivable}, miturut sarat~(1) kita nduweni
$\Th{T} \Proves \OProv[\Th{T}](\gn{!H})$, nanging dudutan kang padha
uga bisa dipikolehi kanthi ngetrapake modus ponens marang ekuivalensi
ing ndhuwur. Mula kita ora oleh dudutan yen $\Th{T}$ ora konsisten,
saora-orane ora kanthi argumen kang padha kaya kanggo ukara G\"odel.
Mesthi wae, iki uga durung nuduhake yen $\Th{T}$ \emph{pancen}
!!{derive}~$!H$.

Kita bisa ngrembakakake pitakon iki sethithik supaya luwih maju.
Arah saka kiwa menyang tengen ing ekuivalensi titik tetep,
$\OProv[\Th{T}](\gn{!H}) \lif !H$, minangka conto saka skema umum
kang diarani \emph{prinsip refleksi}:
$\OProv[\Th{T}](\gn{!A}) \lif !A$.
Jeneng iki digunakake amarga skema kasebut, ing sawijining teges,
nyatakake yen $\Th{T}$ bisa ``ngrefleksikake'' apa kang bisa
!!{derive}; intine, skema iki nyatakake,
``Yen $\Th{T}$ bisa !!{derive}~$!A$, mula~$!A$ bener,''
kanggo saben~$!A$.
Mesthi wae, iki mung bener kanggo teori kang sahih, lan iki menehi
pituduh yen umume teori ora bakal !!{derive} saben contone.
Dadi, sawijining teori (kang cukup kuwat lan nyukupi sarat
!!{derivability}) bisa !!{derive} conto endi?
Mesthi wae kabeh conto kang $!A$-ne dhewe !!{derivable}.
Asil sabanjure nuduhake yen mung kuwi contone.

\begin{thm}\ollabel{thm:lob}
Upamana $\Th{T}$ teori kang !!a{axiomatizable} lan ngembangake
$\Th{Q}$, sarta upamana $\OProv[\Th{T}](y)$ formula kang nyukupi
sarat P1--P3 saka \olref[2in]{sec}. Yen $\Th{T}$ !!{derive}s
$\OProv[\Th{T}](\gn{!A}) \lif !A$, mula sejatine $\Th{T}$
!!{derive}s $!A$.
\end{thm}

Kanthi tembung liya, yen $\Th{T} \Proves/ !A$, mula $\Th{T} \Proves/
\OProv[\Th{T}](\gn{!A}) \lif !A$. Asil iki dikenal minangka
teorema L\"ob.

\begin{explain}
Gambaran intuitif kanggo bukti teorema L\"ob yaiku sawijining bukti
kang cerdik yen Sinterklas ana. (Yen sampeyan ora seneng marang
dudutan iku, sampeyan bisa nggantine nganggo dudutan liya kang
sampeyan karepake.) Buktine mangkene:
\begin{enumerate}
\item Upamana $X$ iku ukara, ``Yen $X$ bener, mula Sinterklas ana.''
\item Upamana $X$ bener.
\item Mula isi ukara kasebut bener; yaiku, yen $X$
  bener, mula Sinterklas ana.
\item Amarga kita nganggep $X$ bener, kanthi modus ponens saka
  (2) lan~(3) kita bisa ndudut yen Sinterklas ana.
\item Kita wis kasil ndeduksi (4), ``Sinterklas ana,'' saka
  anggepan~(2), ``$X$ bener.'' Miturut bukti kondisional,
  kita wis nuduhake: ``Yen $X$ bener, mula Sinterklas ana.''
\item Nanging iki persis ukara~$X$. Dadi kita wis nuduhake yen
  $X$ bener.
\item Mula, miturut argumen (2)--(4) ing ndhuwur, Sinterklas ana.
\end{enumerate}
Yen gagasan iki diformalisasi kanthi ngganti ``bener'' dadi
``!!{derivable},'' lan ``Sinterklas ana'' dadi~$!A$,
kita oleh bukti teorema L\"ob. Carane yaiku ngetrapake lema
titik tetep marang !!{formula}~$\OProv[\Th{T}](y) \lif !A$.
Titik tetep saka formula kasebut cocog karo ukara~$X$ ing
gambaran sadurunge.
\end{explain}

\begin{proof}[Bukti \olref{thm:lob}]
Upamana $!A$ iku !!a{sentence} saengga $\Th{T}$ !!{derive}s
$\OProv[\Th{T}](\gn{!A}) \lif !A$. Upamana $!B(y)$ iku
!!{formula}~$\OProv[\Th{T}](y) \lif !A$, banjur gunakna lema titik
tetep kanggo nemokake !!a{sentence}~$!D$ saengga $\Th{T}$
!!{derive}s $!D \liff !B(\gn{!D})$. Mula saben pratelan ing
ngisor iki !!{derivable} ing $\Th{T}$:
\begin{align}
  & !D \liff (\OProv[\Th{T}](\gn{!D}) \lif !A) \ollabel{L-1}\\
  & \qquad \text{$!D$ iku titik tetep saka~$!B(y)$}\notag \\
  & !D \lif (\OProv[\Th{T}](\gn{!D}) \lif !A) \ollabel{L-2}\\
  & \qquad\text{saka \olref{L-1}}\notag\\
  & \OProv[\Th{T}](\gn{!D \lif (\OProv[\Th{T}](\gn{!D}) \lif !A)}) \ollabel{L-3}\\
  & \qquad \text{saka \olref{L-2} miturut sarat P1}\notag \\
  & \OProv[\Th{T}](\gn{!D}) \lif \OProv[\Th{T}](\gn{\OProv[\Th{T}](\gn{!D}) \lif !A})
  \ollabel{L-4}\\
  &\qquad \text{saka \olref{L-3} kanthi sarat P2}\notag \\
  & \OProv[\Th{T}](\gn{!D}) \lif (\OProv[\Th{T}](\gn{\OProv[\Th{T}](\gn{!D})}) \lif \OProv[\Th{T}](\gn{!A})) \ollabel{L-5}\\
  &\qquad \text{saka \olref{L-4} kanthi P2 maneh} \notag\\
& \OProv[\Th{T}](\gn{!D}) \lif \OProv[\Th{T}](\gn{\OProv[\Th{T}](\gn{!D})}) \ollabel{L-6}\\
  & \qquad\text{miturut sarat !!{derivability} P3} \notag\\
  & \OProv[\Th{T}](\gn{!D}) \lif \OProv[\Th{T}](\gn{!A}) \ollabel{L-7} \\
  &\qquad\text{saka \olref{L-5} lan \olref{L-6}}\notag\\
  & \OProv[\Th{T}](\gn{!A}) \lif !A \ollabel{L-8}\\
  &\qquad\text{miturut anggepan ing teorema} \notag\\
  & \OProv[\Th{T}](\gn{!D}) \lif !A \ollabel{L-9}\\
  &\qquad\text{saka \olref{L-7} lan \olref{L-8}}\notag\\
  & (\OProv[\Th{T}](\gn{!D}) \lif !A) \lif !D \ollabel{L-10}\\
  & \qquad \text{saka \olref{L-1}}\notag \\
  & !D \ollabel{L-11}\\
  & \qquad\text{saka \olref{L-9} lan \olref{L-10}}\notag \\
  & \OProv[\Th{T}](\gn{!D}) \ollabel{L-12}\\
  & \qquad\text{saka \olref{L-11} miturut sarat~P1}\notag \\
  & !A \qquad\qquad\text{saka \olref{L-8} lan \olref{L-12}}\notag
\end{align} Cathetan koreksi sumber OLPL-801: baris pungkasan A ora mung saka L-8 lan L-12; L-7 uga perlu kanggo ngasilake Prov-A saka Prov-D, banjur L-8 menehi A.
\end{proof}

Kanthi teorema L\"ob, kita oleh bukti cekak kanggo teorema ora jangkep
kapindho (kanggo teori kang nduweni predikat !!a{derivability} kang
nyukupi sarat P1--P3): yen $\Th{T} \Proves
\OProv[\Th{T}](\gn{\lfalse}) \lif \lfalse$, mula
$\Th{T} \Proves \lfalse$. Yen $\Th{T}$ konsisten,
$\Th{T} \Proves/ \lfalse$. Mula,
$\Th{T} \Proves/ \OProv[\Th{T}](\gn{\lfalse}) \lif \lfalse$,
yaiku $\Th{T} \Proves/ \OCon[\Th{T}]$. Teorema iki uga bisa
digunakake kanggo nuduhake yen~$!H$, titik tetep saka
$\OProv[\Th{T}](x)$, iku !!{derivable}. Amarga
\begin{align*}
  \Th{T} & \Proves \OProv[\Th{T}](\gn{!H}) \liff !H\\
  \intertext{mligine}
    \Th{T} & \Proves \OProv[\Th{T}](\gn{!H}) \lif !H
\end{align*}
mula miturut teorema L\"ob, $\Th{T} \Proves !H$.

% Going in the other direction, for homework I
% may ask you to work through a short proof of L\"ob's theorem, using
% the second incompleteness theorem instead of the fixed-point lemma.

\begin{prob}
Upamana $\Th{T}$ iku teori kang diaksiomatisasi kanthi komputabel,
lan $\OProv[\Th{T}]$ iku predikat !!a{derivability} kanggo
$\Th{T}$. Gatekna patang pratelan ing ngisor iki:
\begin{enumerate}
\item Yen $T \Proves !A$, mula $T \Proves \OProv[\Th{T}](\gn{!A})$.
\item $T \Proves !A \lif \OProv[\Th{T}](\gn{!A})$.
\item Yen $T \Proves \OProv[\Th{T}](\gn{!A})$, mula $T \Proves !A$.
\item $T \Proves \OProv[\Th{T}](\gn{!A}) \lif !A$
\end{enumerate}
Ing sarat-sarat apa saben pratelan iki bener?
\end{prob}

\end{document}
